\begin{abstract}
Amidst the increasing integration of technology in car interiors, the risk of driver distraction has become a critical concern for automotive safety. Addressing this issue requires robust methods for detecting driver distractions, often relying on intricate models trained on vast amounts of labeled data. 
However, obtaining such labeled data can be expensive and time-consuming.
In this context, self-supervised learning emerges as a promising approach, leveraging unlabeled data to learn meaningful representations and reduce dependency on annotated datasets.
In this study, we explore self-supervised learning methods for 3D skeleton-based driver activity recognition. We evaluate the performance of our proposed SkelDINO-SAM method across diverse backbone architectures.
Utilizing the Drive\&Act dataset, characterized by its long-tailed distribution of activity classes, we evaluate the effectiveness of our approach in addressing challenges associated with real-world scenarios.
Our findings highlight the superiority of transformer-based backbones, particularly when combined with our SkelDINO-SAM approach. 
Through extensive experiments and ablation studies, we demonstrate the efficacy of our method in enhancing driver secondary task recognition accuracy. 
Overall, our results show the efficiency of our approach outperforming the state-of-the-art method by 11.29\%. Our code is published on \href{https://github.com/davidjlerch/darssl}{GitHub}.
\end{abstract}